\section{Introduction}
\label{sec:intro}

Image captioning performance has dramatically improved over the past decade. Despite such impressive results, it is unclear to what extent captioning models actually rely on image content: as we show, existing metrics fall short of fully capturing the captions' relevance to the image. In Figure~\ref{fig:teaser} we show an example where a competitive captioning 
model, Neural Baby Talk (NBT)~\cite{lu2018neural}, incorrectly generates the object ``bench.'' 
%
We refer to this issue as object \emph{hallucination}.
%

While missing salient objects is also a failure mode, captions are summaries and thus generally not expected to  describe all objects in the scene. On the other hand, describing objects that are \emph{not present} in the image has been shown to be less preferable to humans.
%
%
For example, the LSMDC challenge~\cite{lsmdc-2017}
%
documents that correctness is more important to human judges than specificity. 
%
In another study, ~\cite{macleod} analyzed how visually impaired people react to automatic image captions. They found that people vary in their preference of either coverage or correctness. For many visually impaired who value correctness over coverage, hallucination is an obvious concern.
%

%
\begin{figure}[t]
\centering
%
\includegraphics[width=\linewidth]{figures/emnlp18-teaser.pdf}
\caption{\label{fig:teaser} Image captioning models often ``hallucinate'' objects that may appear in a given  context, like e.g. a \emph{bench} here. Moreover, the sentence metrics do not always appropriately penalize such hallucination.  Our proposed metrics (\metric{}s and \metric{}i) reflect hallucination.  For \metric{} \textit{lower is better}.
}
\end{figure}
%
Besides being poorly received by humans, object hallucination reveals an internal issue of a captioning model, such as not learning a very good representation of the visual scene or overfitting to its loss function. 

%
In this paper we assess the phenomenon of object hallucination in contemporary captioning models, and consider several key questions.
The first question we aim to answer is: \emph{Which models are more prone to hallucination?} We analyze this question on a diverse set of captioning models, spanning different architectures and learning objectives. To measure object hallucination, we propose a new metric, \emph{\metric{} (\metricfullname{})}, which captures image relevance of the generated captions. Specifically, we consider both ground truth object annotations (MSCOCO Object segmentation~\cite{lin2014microsoft}) and ground truth sentence annotations (MSCOCO Captions~\cite{chen2015microsoft}).
%
Interestingly, we find that models which score best on standard sentence metrics do not always hallucinate less.%

The second question we raise is: \emph{What are the likely causes of hallucination?} %
While hallucination may occur due to a number of reasons, we believe the top factors include visual misclassification and over-reliance on language priors. The latter may result in memorizing which words ``go together'' regardless of image content, which may lead to poor generalization, once the test distribution is changed. We propose \emph{image and language model consistency} scores to investigate this issue, %
and find that models which hallucinate more tend to make mistakes consistent with a language model.%

Finally, we ask: \emph{How well do the standard metrics capture hallucination?} It is a common practice to rely on automatic sentence metrics, e.g. CIDEr \cite{cider}, to evaluate  captioning performance during development, and few employ human evaluation to measure the final performance of their models. As we largely rely on these metrics, it is important to understand how well they capture the hallucination phenomenon. In Figure \ref{fig:teaser} we show how two sentences, from NBT with hallucination and from TopDown model~\cite{anderson2017bottom} -- without, are scored by the standard metrics. As we see, hallucination is not always appropriately penalized. We find that by using additional ground truth data about the image in the form of object labels, our metric \metric{} allows us to catch discrepancies that the standard captioning metrics cannot fully capture. We then investigate ways to assess object hallucination risk with the standard metrics. Finally, we show that \metric{} is complementary to the standard metrics in terms of capturing human preference.

%
%
%
%

%
%
%
%
%
%
%
%
%

%
%
%
%
%


%

%
%

%
%
%
%
%
%
%

%
%

%
%
%
%
